\section{Planar routing and invariant coding planes}\label{sec:sheets-planes}
The contraction, parking and expansion strategy has an area-preserving predecessor in \cite[Proposition~5.1]{CardonaMirandaPeraltaSalasPresas2021}. The proof below gives the needed planar motion explicitly; the three-dimensional incompressibility is supplied separately by the invariant-plane lift.
\begin{lemma}[Planar routing through parking positions]\label{lem:sa-planar-routing}
Let $D$ be an open convex rational rectangle.  Let $B_i$ and $C_i$ be
two separately disjoint finite families of nondegenerate rational closed
rectangles compactly contained in $D$.  An effective smooth planar field,
supported in $(0,1)\times D$, realizes their matched positive diagonal
affine maps on positive neighborhoods of all $B_i$.  If $G\subset D$ is
an open convex rational rectangle, every pair $B_i,C_i\Subset G$ can have
its entire controlled motion confined to $G$.
\end{lemma}
\begin{proof}
Take $G=D$ if no smaller safe rectangle is specified.
For an empty family use zero.  Otherwise choose distinct rational
parking points in $G$ avoiding all source and
target centers.  Move centers first from the sources to their parking
points, one at a time, and then from parking to their targets, one at a
time.  Each destination is unoccupied.  For a move $p\to q$, a two-segment
path $p\to a\to q$ avoids the finite set of stationary centers if $a$
avoids the lines joining an endpoint to one of those centers.  Choose
$a$ rational in $G$ when both endpoints are there, and in $D$ otherwise.
To make this choice effective, take a rational parabola segment in the
chosen open rectangle: every forbidden line meets it at at most two
points, so finitely many rational samples suffice.  Convexity keeps the
segments inside the chosen rectangle.

All segment-to-stationary-center squared distances and boundary
clearances have positive rational lower bounds.  Choose a rational
$\epsilon>0$ smaller than every endpoint half-side, with $10\epsilon$
below every relevant distance and boundary clearance.  Shrink each
source rectangle to the square of half-side $\epsilon$ about its center,
translate these squares along the chosen segments in order, and expand
inside the target rectangles.  During contraction the source cutoffs have pairwise disjoint supports;
during expansion the target cutoffs do.  The moving square's support
of half-side $2\epsilon$ misses each stationary square.  During a shrink/expansion about $c$ with
positive half-widths $v_a(t)$, use a cutoff times
\[
 \left(\frac{\dot v_1}{v_1}(x_1-c_1),
       \frac{\dot v_2}{v_2}(x_2-c_2)\right);
\]
during a translation use a moving plateau times $\dot c$.
The exact path is $c(t)+\diag(v_1(t),v_2(t))y$, $y\in[-1,1]^2$.
Its endpoint factors are the ratios of the target to source half-sides.
Flat schedules on finitely many successive slots make a smooth pulse.
Shrinking and expanding stay in the endpoint rectangles; the prescribed
path clearances prove the extra confinement in $G$.

Widths may be interpolated exponentially or linearly.  With padded
plateaus, either choice gives an effective positive initial neighborhood
by the finite products of affine norms.  A cutoff that equals one only
on the closed moving rectangle still proves the same formula there,
including its boundary, but by itself asserts no exterior neighborhood.
We shall state explicitly when such a closed-rectangle implementation
is being retained.  The empty-family case uses the zero field.
\end{proof}

\begin{lemma}[Invariant-plane lift]\label{lem:sa-invariant-lift}
Let $b(t,r)$ be an effective smooth planar field, periodic on $\mathbb T^2$
or compactly supported in $\mathbb R^2$, with bounded mixed derivatives.
Let $h$ be smooth effective with bounded derivatives, $h(z_0)=0$ and
$h'(z_0)=1$.  Then
\begin{equation}\label{eq:sa-lift}
 U(t,r,z)=(h'(z)b_1,h'(z)b_2,-h(z)\operatorname{div}_r b)
\end{equation}
is divergence free and carries exactly the planar flow on $z=z_0$.
The choice $h(z)=\sin(2\pi z)/(2\pi)$ for $z_0=0$ gives a mean-zero
unit-torus field.  The choice $h=(z-z_0)\chi(z)$, with $\chi=1$ near
$z_0$ and compact support, gives compact support on $\mathbb R^3$ when
$b$ is compactly supported, or a mean-zero torus field when $h$ is
supported in one vertical chart.  Time periods and mixed-derivative
bounds are preserved.
\end{lemma}
\begin{proof}
The divergence is $h'\operatorname{div}b-h'\operatorname{div}b=0$.
On the indicated plane $U=(b,0)$, so ODE uniqueness gives its invariance.
For the sine choice both vertical factors integrate to zero.  For compact
$h$, $\int h'=0$ and $\int\operatorname{div}b=0$ give the three means.
All remaining assertions follow by differentiation of finite products.
\end{proof}

\section{Histories above and below separators}
\label{sec:sa-separator-histories}
A separator lets one stack contain both work symbols and a finite history
without mixing their roles.  Storing history below the separator makes
the scan length depend on the updated work word.  Storing it above the
separator makes the scan traverse the accumulated history instead.  This
choice changes the instruction count and the unused image cylinders
available for initialization.

\subsection{History below a delimiter}\label{subsec:sa-delimiter}
A delimiter lets a finite left work word sit above its history.  A marker
on the right remembers how far the subsequent return scan must go.
Use work alphabet $\Gamma$ and normalize to one halt state $q_h$.
Introduce $\#,\star$ and records
$h_i$.  For a right or stay instruction $\delta(q,a)=(q',b,d)$ let
$i=(q,a)$; for a left instruction let $i=(q,a,c)$ with
$c\in\Gamma\sqcup\{\#\}$.  Write $q^+(i)=q'$.

\begin{lemma}[Delimiter-history simulation]\label{lem:sa-delimiter}
Starting at $(N_{q_0},\#\square^\infty,w\square^\infty)$, normal
configurations have form $(N_q,L\#H\square^\infty,R)$, where $L,H$
are finite, $R$ is eventually blank, and the represented left tape is
$L\square^\infty$.  One machine instruction takes exactly $2|L'|+3$
rules, where $L'$ is the updated finite left work word.  Full source
and image cylinders are separately disjoint, and the initialized run
reaches $N_{q_h}$ exactly when the machine halts.
\end{lemma}
\begin{proof}
For $R=aR_*$ the first rule gives the updated pair
\[
 (L',R')=\begin{cases}
 (bL,R_*),&d=+,\\
 (L,bR_*),&d=0,\\
 (L_*,cbR_*),&d=-,\ L=cL_*,\\
 (\epsilon,\square bR_*),&d=-,\ L=\epsilon.
 \end{cases}
\]
It also places $\star$ before $R'$.  The $D_i$ scan transfers $|L'|$
work symbols until $\#$ is exposed, inserts $h_i$ below it, and enters
$U_{q^+(i)}$.  Exactly $|L'|$ reverse transfers restore $L'$, and the
marker-removal rule finishes the step.  This proves the count and tape
invariant, including the new blank at the left boundary.  Record
displacements recover the absolute head position.

The complete rules are
\begin{equation}\label{eq:sa-delimiter-table}
\begin{aligned}
 (N_q;\epsilon,a)&\to(D_i;b,\star),&&d=+,\\
 (N_q;\epsilon,a)&\to(D_i;\epsilon,\star b),&&d=0,\\
 (N_q;c,a)&\to(D_i;\epsilon,\star cb),&&d=-,\ c\in\Gamma,\\
 (N_q;\#,a)&\to(D_i;\#,\star\square b),&&d=-,\ c=\#,\\
 (D_i;c,\epsilon)&\to(D_i;\epsilon,c),&&c\in\Gamma,\\
 (D_i;\#,\epsilon)&\to(U_{q^+(i)};\#h_i,\epsilon),\\
 (U_q;\epsilon,c)&\to(U_q;c,\epsilon),&&c\in\Gamma,\\
 (U_q;\epsilon,\star)&\to(N_q;\epsilon,\epsilon).
\end{aligned}
\end{equation}
No normal rule leaves $N_{q_h}$.

The tested top symbols separate sources, with an extra left-top test
for left moves.  Into $D_i$ a normal entry starts with $\star$ on the
right and a loop with an ordinary letter.  Into $U_q$ a history entry
starts with $\#h_i$ on the left and a loop with an ordinary letter.
Distinct records identify distinct entries.  Into $N_q$ there is just
its marker-removal rule.  These are inverse checks for arbitrary tails,
so induction gives all assertions including initial halting.
\end{proof}

Three useful choices of observation follow from this same delimiter
mechanism.  One uses its table directly; a second changes its coordinates;
a third separates the work update and marker push into two genuine rules.
This last modification changes the sample times even though contracting
those two rules recovers the original table.

\begin{theorem}[Delimiter processors on invariant planes]
\label{thm:sa-delimiter}\label{thm:sa-delimiter-origin}
\label{thm:sa-left-stack-plane}
For every machine and finite input, each row below gives an effective
smooth mean-zero general force on the unit torus, periodic of period one
for $t\ge1$, with all mixed derivatives bounded.  Its zero-data solution
is unique in the classical periodic class and has uniformly bounded
energy.  The stated fixed particle enters the stated fixed open set
exactly when the machine halts.
\begin{center}\small
\begin{tabular}{lll}
Processor & Initial label & Open set\\\hline
Delimiter & $(1/2,1/2,0)$ & $(2/3,5/6)^2\times\mathbb T$\\
Delimiter, origin coordinates & $(0,0,0)$ & $\{x:2/3<x_1<11/12\}$\\
Separate marker push & $(1/4,1/4,0)$ & $\{x:3/4<x_1<1\}$
\end{tabular}\end{center}
The repeated pulse uses only the machine; a separate loader uses the
input.  All formulas allow arbitrary fixed positive real viscosity as
a parameter; absolute computability uses computable viscosity.
\end{theorem}
\begin{proof}
\emph{The third processor.}  For a branch $i$ of
\eqref{eq:sa-delimiter-table}, replace its normal rule by a rule that
performs the same tape update but enters a fresh state $B_i$ without
pushing $\star$, followed by
$(B_i;\epsilon,\epsilon)\to(D_i;\epsilon,\star)$.
Use exactly the other rows of that table.  Every $B_i$ has one incoming
and one outgoing rule.  Into $D_i$ its new entry has right top $\star$,
whereas loop entries have ordinary right top, so the previous full-cylinder
inverse proof remains valid.  On every full source cylinder, the composite
of the two new rules equals the old normal rule.  This proves the exact
relationship, and the machine-step count is now $2|L'|+4$.


\emph{Codes and placement.}  Use odd digits and base $B=2m+1$, where
$m$ is the relevant full alphabet size.  Apply Lemma~\ref{lem:sa-radix}
in the following state rectangles.
For the first row take the halt square $[7/10,4/5]^2$.  For the $M$
other controls take centers $(1/8+(2j+1)/(8M),1/4)$,
$0\le j<M$, and half-side $1/(16M)$.
For the second row take the halt rectangle
$[3/4,7/8]\times[1/4,1/2]$ and other rectangles
\[
 [1/8+j/(4M),1/8+j/(4M)+1/(8M)]\times[1/4,1/2].
\]
For the third row give the halt control the square centered at
$(7/8,1/2)$ of side $1/16$, and the other $M$ controls squares centered
at $(i/[2(M+1)],1/2)$, $1\le i\le M$, of side $1/[8(M+1)]$.
The first two nonhalt families lie in $(0,1/2)^2$, respectively
$(0,1/2)\times(0,1)$; the third lies in the latter open rectangle.
In each case the full source and target rectangle families are separately
separated by the proved cylinder checks and digit gaps.

\emph{Routing and loading.}  Apply Lemma~\ref{lem:sa-planar-routing}
in $(0,1)^2$, using the corresponding nonhalt open rectangle as the safe
region.  This proves exact maps on full rectangles and all-time confinement
whenever both endpoints are nonhalting.  Both its exponential-width and
linear-width implementations are available.  The second row also permits
the following explicit linear-width choice: park strictly to the left of
all endpoint centers and choose every bend to the left of both endpoints,
off the finite forbidden lines.  Nonhalt centers then remain below $3/8$;
small squares of half-side below $1/8$ stay below $1/2$.  This gives the
same endpoint maps with a directly checked closed-rectangle plateau.
For the third row one may choose rational bends on $(\lambda s,s^2)$,
$\lambda=1/2$ or $1$.  After clearing denominators, forbidden lines give
nonzero integer polynomials of degree at most two.  A rational $s=1/D$
with $D$ larger than all absolute leading coefficients cannot be a root
by the rational-root theorem.  Avoid the finite endpoint denominators
as well and take $D$ large enough for the prescribed open rectangle.
This is an effective form of the parking choice, with no generic-position
assumption.

All initial codes are rational.  In the first row a planar compact
translation, with plateau half-side $1/100$ and support half-side $2/100$,
loads from $(1/2,1/2)$.  In the second row the spatially constant planar
velocity $\dot\theta(t)r_{\rm in}$ loads from the origin.  In the third
row a small compact translation tube loads from $(1/4,1/4)$.
The loader paths avoid the respective observers when the initial state
is nonhalting.  Follow the loader by the repeated unit pulse and apply
Lemma~\ref{lem:sa-invariant-lift} with $h=\sin(2\pi z)/(2\pi)$.
The three velocities are mean zero, smooth at all joins, and zero
initially; Lemma~\ref{lem:sa-realization} supplies their forces.

At time $1+k$ the particle is the code after $k$ primitive rules until
halting.  A halt code lies strictly in the corresponding observed set,
including an initial halt at time one.  For nonhalting runs, the loader
and every subsequent controlled path avoid that set by the established
safe-region bounds.  This proves the continuous-time equivalences and
all three force contracts.
\end{proof}

\subsection{An input instruction inside the repeating program}
\label{subsec:sa-prefix-insertion}
A fixed initial code can instead load the input through an ordinary
prefix rule.  This makes the whole force periodic from zero, at the cost
of placing the input in its repeating finite table.
Use a new initial state $S$, and the delimiter alphabet renamed
$\sharp,\diamond,h_j$.  The input rule is
\[
 (S;\epsilon,\epsilon)\to(F_{j_*};\epsilon,\diamond w),
 \qquad q(j_*)=q_0.
\]
For each nonhalt instruction $\delta(q,a)=(q',b,d)$ use primary rules
\[
\begin{array}{ll}
 d=+:&(q;\epsilon,a)\to(F_j;b,\diamond),\\
 d=0:&(q;\epsilon,a)\to(F_j;\epsilon,\diamond b),\\
 d=-:&(q;c,a)\to(F_j;\epsilon,\diamond cb),\quad c\in\Gamma,\\
     &(q;\sharp,a)\to(F_j;\sharp,\diamond\square b).
\end{array}
\]
Every alternative, including initialization, has its own $j$ and history
letter $h_j$, and $q(j)=q'$.  Complete the table by
\begin{equation}\label{eq:sa-prefix-insertion-scans}
\begin{array}{ll}
 (F_j;c,\epsilon)\to(F_j;\epsilon,c),&c\in\Gamma,\\
 (F_j;\sharp,\epsilon)\to(G_{q(j)};\sharp h_j,\epsilon),\\
 (G_q;\epsilon,c)\to(G_q;c,\epsilon),&c\in\Gamma,\\
 (G_q;\epsilon,\diamond)\to(q;\epsilon,\epsilon).
\end{array}
\end{equation}
There are no rules from halt controls or into $S$.

\begin{lemma}\label{lem:sa-prefix-insertion}
These source and image cylinders are separately disjoint.  From
$(S,\sharp\square^\infty,\square^\infty)$ the first procedure inserts
$w$, and thereafter one ordinary step takes $2|L'|+3$ rules and reaches
the correct main configuration.  The run reaches a halt main control
exactly when the machine halts.
\end{lemma}
\begin{proof}
Every $F_j$ has a unique primary entry with right top $\diamond$;
its loops have ordinary right tops.  Entries into $G_q$ have distinct
left prefixes $\sharp h_j$, or ordinary left tops for loops.  Each
main state has its unique exit from $G_q$.  Source separation is by the
same tested top symbols as in Lemma~\ref{lem:sa-delimiter}, with the
single additional source state $S$.  These checks prove the full-domain
claim.  At main states the stacks are $(L\sharp H\square^\infty,R)$.
The primary update gives $(L'\sharp H\square^\infty,\diamond R')$;
the two scans transfer and restore $L'$, insert $h_j$, and delete the
marker.  Their count and interpretation are those proved for the
delimiter table.  Initialization is its case $L'=\epsilon$,
$R'=w\square^\infty$ and takes three rules.  The history records include
the zero-displacement initialization and recover the absolute head.
\end{proof}

\begin{theorem}[Input insertion in a periodic compact force]
\label{thm:sa-prefix-insertion}
This processor gives an effective smooth general force on $\mathbb R^3$
with period one from zero, a fixed compact spatial support, and all mixed
derivatives bounded.  Its zero-data solution is globally smooth, has
uniformly bounded energy, and is unique among smooth classical solutions
with $u\in C_tH^2\cap C_t^1L^2$, $p\in C_tH^1$, and bounded velocity
and gradient on each finite time interval.  The particle from the origin has
first coordinate greater than $2$ at some time exactly when the machine
halts.  The formula is effective relative to any fixed real $\nu>0$.
\end{theorem}
\begin{proof}
For $m$ letters use digits $2,4,\ldots,2m$ and base $K=2m+2$.
Set $a_S=-E(\sharp\square^\infty)$, give halt main states distinct
offsets $4,8,\ldots$, and other states offsets $-4,-8,\ldots$.
With $b_0=-E(\square^\infty)$, encode by
$(a_s+E(L),b_0+E(R),0)$, so the fixed initial code is zero.
Lemma~\ref{lem:sa-radix} gives separated source and target rectangles.
Thicken to $R_i^-\times[-1,1]$ and apply
Lemma~\ref{lem:sa-full-box-layers} with scales
$K^{|v|-|v'|},K^{|w|-|w'|}$ and their reciprocal product.
For example take the three switches $\sigma(8t-(2j-1))$, $j=1,2,3$,
and collar one quarter of the minimum of $1$ and the rectangle gaps.
The field is zero for $t\le1/8$ and $t\ge6/8$ within a period.
Its periodization is smooth from zero, so
Lemma~\ref{lem:sa-realization} applies.

Induction at integer times follows the initialized rule sequence.
A halt code has first coordinate at least $4$.  Every nonhalting used
branch has endpoint centers at most $1$ and widths at most $1$;
the exponential-width bound gives first coordinate at most $3/2$
throughout.  The force therefore has the required all-time event.
The complete input-insertion branch is included in every period, and
its actual one-time use follows from the symbolic dynamics.
\end{proof}

\subsection{A finite work window on each side}\label{subsec:sa-delimiter-history}
The preceding delimiter processor leaves the right tape infinite.  If both
work words are finite, crossing either end must instead insert a new blank.
This changes the guards of the right-move rules and introduces a distinct
marker-push state.  Here the work alphabet is $\Gamma$ and the sole halt
state is denoted by $h$.

At a ready control $A_q$ use stacks
$L\#H\square^\infty$ and $R\#\square^\infty$, where $L,H$ are finite
and $R$ is a nonempty work word.  Initially $L=H=\epsilon$ and $R=w$,
except that empty input is replaced by the one-letter word $\square$.
Denote this initial right work word by $R_{\rm in}$.
For $\delta(q,a)=(q',b,d)$ the work prefix replacements are
\[
\begin{array}{c|c|c|l}
 d&(v,w)&(v',w')&\text{condition}\\\hline
 0&(\epsilon,a)&(\epsilon,b)&\\
 +&(\epsilon,ac)&(b,c)&c\in\Gamma\\
 +&(\epsilon,a\#)&(b,\square\#)&\\
 -&(c,a)&(\epsilon,cb)&c\in\Gamma\\
 -&(\#,a)&(\#,\square b)&
\end{array}
\]
Give every alternative its own index $r$, control $P_r$, and record
$\rho_r$.  Its state change is $A_q\to P_r$ and $q'(r)=q'$.
Add a marker $J$ and controls $S_r,T_q$, where $q$ runs through every
machine state, including $h$.  The rules are
\[
\begin{array}{ll}
 (P_r;\epsilon,\epsilon)\to(S_r;\epsilon,J),\\
 (S_r;c,\epsilon)\to(S_r;\epsilon,c),&c\in\Gamma,\\
 (S_r;\#,\epsilon)\to(T_{q'(r)};\#\rho_r,\epsilon),\\
 (T_q;\epsilon,c)\to(T_q;c,\epsilon),&c\in\Gamma,\\
 (T_q;\epsilon,J)\to(A_q;\epsilon,\epsilon).
\end{array}
\]
The full alphabet is $\Gamma_{\rm full}=\Gamma\sqcup\{\#,J\}
\sqcup\{\rho_r\}_r$, with all added symbols distinct.
There is no work rule from $A_h$.
The work rule updates the finite tape interval, adjoining a blank exactly
when a move crosses one of its ends.  After pushing $J$, the $S_r$ scan
transfers the updated left word, inserts $\rho_r$ below its delimiter,
and the $T_q$ scan restores it and removes $J$.  Hence one step takes
$4+2|L'|$ rules and leaves $R'$ nonempty.  All finite scans terminate.
Sources at $A_q$ are separated by the scanned symbol and the displayed
additional boundary test.  Each $P_r$ has one incoming rule.  Images
into $S_r$ have right top $J$ or an ordinary letter; images into $T_q$
have left prefix $\#\rho_r$ or an ordinary letter.  Each $A_q$ has one
marker-removal entry.  These observations prove source and image
separation for arbitrary tails, as well as the initialized halt equivalence.

\begin{theorem}[Finite windows and a fixed slab]\label{thm:sa-delimiter-history}
This processor gives an effective smooth mean-zero force on the unit
torus, periodic of period one for $t\ge1$, with every mixed derivative
bounded.  Its zero-data solution is globally smooth, unique in the
classical periodic class, and has uniformly bounded energy.  The particle
from $(1/4,1/4,1/4)$ enters $\{x:1/2<x_1<3/4\}$ exactly when the
machine halts.  A general residual force with pressure zero or a
solenoidal projected force with its accompanying pressure may be chosen;
both give the same velocity.  The repeated pulse depends only on the
table and the loader on the input.
\end{theorem}
\begin{proof}
Let $K$ be the number of controls, $\eta=1/(16K)$, and use odd digits
with base $B=2|\Gamma_{\rm full}|+1$.  Give nonhalt controls offsets
$a_{s_j}=1/4+2j\eta$ and set $a_{A_h}=5/8$.  The code is
\[
 (a_s+\eta E(L),1/4+\eta E(R),1/4).
\]
All nonhalt state intervals are in $[1/4,3/8)$, while the halt interval
lies in $(1/2,3/4)$.  Lemma~\ref{lem:sa-radix} gives separately separated
rectangles $E_i^0,E_i^1\subset[1/4,3/4]^2$ and scale factors
$\alpha_i=B^{|v_i^0|-|v_i^1|}$,
$\beta_i=B^{|w_i^0|-|w_i^1|}$.  Write $F_i$ for the positive
diagonal affine map $E_i^0\to E_i^1$.
Here is an explicit thin-box instance of normal compensation.  If there
are $m$ rules, put
\[
 Z_i=\tfrac12+\frac{i}{4(m+1)},\quad
 M=\max_i\max(1,1/\alpha_i)\max(1,1/\beta_i),\quad
 \delta_z=\frac1{64(m+1)M},
\]
\[
 \mu=\min\{d_*/4,1/[64(m+1)]\},
\]
where $d_*$ is the minimum of $1$ and the within-family planar
sup-norm distances.  Lift the boxes
$E_i^0\times[1/4-\delta_z,1/4+\delta_z]$ to $Z_i$, use the horizontal
scales $1-s+s\alpha_i$, $1-s+s\beta_i$ and reciprocal vertical scale,
and lower.  Each vertical half-size is at most $\delta_zM$.
During ascent/descent use the endpoint rectangle gaps; in the middle
stage use the height gaps $1/[4(m+1)]$.  Enlarging by $\mu$ preserves
separation and stays in the cube.  Thus Lemma~\ref{lem:sa-curl-motion}
gives exact affine motion on the full boxes and positive collars.
On their central sheet the first coordinate is
$(1-s)X_1+sF_i(X)_1$.

The input point is the code of
$(A_{q_0},\#\square^\infty,R_{\rm in}\#\square^\infty)$.
A curl translation with plateau half-side $1/32$ and margin $1/32$
loads it from the fixed label during $[0,1]$.  All paths stay inside
the cube.  Repetition gives a mean-zero velocity and
Lemma~\ref{lem:sa-realization} gives its residual force $r$ and pressure
zero.  For the second choice, Proposition~\ref{prop:projection-l2}
solves $\Delta\psi=\operatorname{div}r$, $\int\psi=0$, and gives
\[
 f_{\rm sol}=r-\nabla\psi,\qquad p=-\psi.
\]
It preserves all stated bounds and eventual periodicity, with precisely
this change in pressure.

Induction at times $1+k$ gives the prefix run until halt.  A halt is in
the slab, including an initial halt.  For a nonhalting run the loader
stays below $1/2$ and each subsequent central-sheet coordinate remains
between two nonhalt endpoints.  No intermediate time enters the slab.
The history recovers head displacement; the finite words above the two
delimiters recover the full represented tape.
\end{proof}

\subsection{History above a separator and a repeated loader}
\label{subsec:sa-repeated-loader}
If history is above the left tape separator, every instruction scans its
entire current length.  The number of primitive steps is then explicit.
This also provides a free target cylinder in which a loading branch can
be included in every period.

Use work alphabet $\Gamma$ and normalize to one halt state $q_H$.
Write the ordinary instructions as
$j:(r_j,b_j^{\rm in})\mapsto(r_j^+,b_j^{\rm out},d_j)$.
Adjoin a history alphabet $G=\{g_j\}_j$ and markers $M_L,M_R$,
all disjoint from $\Gamma$.
Let $I_j=\{\perp\}$ for $d_j\ne-1$, and $I_j=\Gamma$ for $d_j=-1$.
Besides the main controls use $D_j,E_{j,c}$.
Here $S_{j,c}$ is $\epsilon$, $b_j^{\rm out}$, or $cb_j^{\rm out}$
for moves $1,0,-1$, respectively.

At a main visit the stacks are $(HM_LL,R)$: $L$ is the infinite reversed
left tape, $R$ starts at the head, and $H$ is the finite history, newest
first.  Entry removes the scanned symbol and inserts $M_R$ on the right.
The $D_j$ scan transfers $H$, then changes the left tape by pushing
$b_j^{\rm out}$, leaving it alone, or popping $c$, according to the move.
The $E_{j,c}$ scan restores $H$ and its exit gives
$(g_jHM_LL',S_{j,c}R_*)$.  This is exactly the next tape configuration
and costs $2|H|+3$ rules.  Initially $H=\epsilon$,
$L=\square^\infty$, $R=w\square^\infty$.

The complete rules are
\[
\begin{aligned}
 (r_j;\epsilon,b_j^{\rm in})&\to(D_j;\epsilon,M_R),\\
 (D_j;g,\epsilon)&\to(D_j;\epsilon,g),&&g\in G,\\
 (D_j;M_L,\epsilon)&\to(E_{j,\perp};M_Lb_j^{\rm out},\epsilon),&&d_j=1,\\
 (D_j;M_L,\epsilon)&\to(E_{j,\perp};M_L,\epsilon),&&d_j=0,\\
 (D_j;M_Lc,\epsilon)&\to(E_{j,c};M_L,\epsilon),&&d_j=-1,\ c\in\Gamma,\\
 (E_{j,c};\epsilon,g)&\to(E_{j,c};g,\epsilon),&&g\in G,\\
 (E_{j,c};\epsilon,M_R)&\to(r_j^+;g_j,S_{j,c}),&&c\in I_j,
\end{aligned}
\]
No rule leaves $q_H$.

Sources are separated by the ordinary read symbol, history tops, markers,
and the extra left-neighbor test.  For images into $D_j$, the right top
is $M_R$ after entry and $g$ after a loop.  Into $E_{j,c}$ the left top
is $M_L$ after modification and $g$ after restoration.  Into a main
control the new left record $g_j$ identifies the instruction; for a
left-moving instruction the first right symbol further identifies $c$.
These checks prove full source and image separation.  They also give the
initialized tape invariant and its exact halt equivalence.

\begin{theorem}[A periodic processor with quadratic sample times]
\label{thm:sa-repeated-loader}
Every machine and finite input give an effective smooth mean-zero general
force on the unit torus, periodic of period one from time zero.  Force
and velocity have all mixed derivatives bounded and vanish near integer
times.  Their spatial supports lie in one compact subset of the cube
chart.  The unique classical periodic zero-data solution is globally
smooth with uniformly bounded energy.  From the fixed label
$(1/2,1/2,1/4)$, the event $3/4<X_1(t)<7/8$ occurs exactly when the
machine halts.  Whenever reached, the configuration after $k$ ordinary
steps is encoded at time $(k+1)^2$.
\end{theorem}
\begin{proof}
For the full alphabet of size $m$ use odd digits and base $B=2m+1$.
If there are $m_s$ controls, put $s_0=1/[16(m_s+1)]$.
The halt square has corner $(13/16,1/4)$ and side $s_0$; all other
squares have corners $(1/8+2js_0,1/4)$ in their enumeration.  Denote
each assigned corner by $p_s$. The nonhalting corners lie
below first coordinate $1/4$.  Encode at height $z_0=1/4$ by
\[
 (p_s+s_0(E(A),E(B)),z_0).
\]
Lemma~\ref{lem:sa-radix} supplies all rule rectangles and their exact
positive diagonal maps.

Add a pair of equal-size loader squares centered at $(1/2,1/2)$ and
\[
 y_{\rm in}=p_{q_0}+s_0
       (E(M_L\square^\infty),E(w\square^\infty)).
\]
Its common half-side is half the minimum of $1/16$ and the distances
from $y_{\rm in}$ to the sides of
$R_{\rm sep}=p_{q_0}+s_0(I(M_L)\times[0,1])$.
These distances are positive, since infinite odd-digit codes lie strictly
inside their prefix intervals.  The loader source misses every state
square.  Its target lies in $R_{\rm sep}$, whereas all ordinary targets
in $q_0$ have a history letter first on the left.  Digit gaps therefore
keep this new target separate.  This verifies injectivity of the enlarged
finite geometric family; it is the reason the loader may repeat.

Apply Lemma~\ref{lem:sa-sheet-routing} to this family, including the
loader.  More explicitly, if there are $N$ pairs, choose heights
$h_i=1/2+i/[4(N+1)]$.  Write $F_i$ for the affine map of pair $i$.  For its horizontal scale
factors $\lambda_{ia}$ let
\[
 \Lambda=\max(\{1\}\cup\{\lambda_{ia}:1\le i\le N,\ a=1,2\}),\quad
 \Lambda_z=\max_i[\min(1,\lambda_{i1})\min(1,\lambda_{i2})]^{-1}.
\]
If $g$ is the minimum of $1$, horizontal boundary clearances, and the
positive within-family coordinate-separating gaps, take horizontal
source padding $g/(10\Lambda)$ and vertical half-thickness
$1/[40(N+1)\Lambda_z]$.
The linear horizontal scale $1-s+s\lambda_{ia}$ and reciprocal normal
scale keep endpoint padding at most $g/10$ and middle-stage half-thickness
at most one tenth of the height gap.  The moving supports are consequently
separate inside the cube.  Use ascent, transformation, descent switches
on $[1/8,2/8]$, $[3/8,5/8]$, $[6/8,7/8]$.
The curl formula gives a mean-zero periodic velocity with zero collars
at integer times.  Its central-sheet horizontal path is
\begin{equation}\label{eq:sa-loader-convex-path}
       (1-s)y+sF_i(y).
\end{equation}
Lemma~\ref{lem:sa-realization} supplies the force, with pressure zero.

The first period loads the input.  Thereafter each period performs the
next primitive rule.  After $r$ ordinary steps $|H|=r$, so the sample
time after $k$ steps is
$1+\sum_{r=0}^{k-1}(2r+3)=(k+1)^2$.
A halt code lies in the observed arc, also when $k=0$.
If the machine does not halt, all endpoints after loading have first
coordinate below $1/4$, while the initial one is $1/2$.
Equation~\eqref{eq:sa-loader-convex-path} excludes the observed arc at
every intermediate time.  History displacements recover the absolute
head index, and removing $H,M_L$ recovers the tape.  All branch and
loader data were prepared from finite input, with no run computed.
\end{proof}

\section{History tracks and return scans}
\label{sec:sa-track-histories}
A separate tape track can store history while the work symbols remain in
their indexed cells.  Both processors in this section mark the work head,
travel to a rightward history frontier, and return before completing the
original head move.  One advances the frontier through a separate control;
the other edits its neighboring cell directly, which requires a different
inverse guard on the return scan.

\subsection{A marked head and a separate log track}
\label{subsec:sa-marked-head-plane}
The first compiler uses a separate instruction to advance the history
frontier before returning to the marked work cell.

Let $A$ be the work alphabet, $Q$ the states, $H$ the halt set, and
$D=\{-1,0,1\}$.  For $e=(q,s,a)\in(Q\setminus H)\times D\times A$
write $(r_e,b_e,d_e)=\delta(q,a)$.  The full tape alphabet is
$\Sigma=A\times(\{\ell,\#\}\sqcup\{[e]\})\times\{0,1\}$,
where the last component marks the return location.
Use controls $S_{q,s},G_e,T_{r,d},L_{r,d}$.

Initialize at $S_{q_0,0}$, head zero, with the prescribed work tape,
all marks zero, and log blank except for $\#$ at cell one.
At the $k$th main visit, records occupy $1,\ldots,k$, the frontier is
$j=k+1$, and the ordinary head satisfies $|i|\le k$.  Thus $i<j$.
The first rule writes the work symbol and marks $i$, $G_e$ reaches the
frontier, $T$ advances it by one, and $L$ returns to the sole mark,
clears it, and performs the original move $d_e$.  Every scanned symbol
satisfies its displayed guard.  These finite scans establish the next
main configuration and the invariant by induction.  Halting is exactly
arrival at an $S_{q,s}$ with $q\in H$; this includes an initial halt.

The following instructions have output order (control, written symbol, move):
\[
\begin{array}{rcll}
 (S_{q,s},(a,h,0))&\mapsto&(G_e,(b_e,h,1),1)&h\ne\#,\\
 (G_e,(c,h,0))&\mapsto&(G_e,(c,h,0),1)&h\ne\#,\\
 (G_e,(c,\#,0))&\mapsto&(T_{r_e,d_e},(c,[e],0),1),\\
 (T_{r,d},(c,\ell,0))&\mapsto&(L_{r,d},(c,\#,0),-1),\\
 (L_{r,d},(c,h,0))&\mapsto&(L_{r,d},(c,h,0),-1)&h\ne\#,\\
 (L_{r,d},(c,h,1))&\mapsto&(S_{r,d},(c,h,0),d)&h\ne\#.
\end{array}
\]
The first row is restricted to nonhalt $q$, and no other rules exist.

All incoming moves to $G_e,T_{r,d}$ are rightward, those to $L_{r,d}$
leftward, and those to $S_{r,d}$ equal $d$.  Moreover target control and
written whole symbol recover the source rule.  At $G_e$, the entering
rule writes mark one while its loops write zero, and $e$ recovers the
old work symbol and control.  At $T$ the record recovers $e$.  At $L$
an entry writes $\#$ and a loop writes a different history symbol.
At $S$ the sole type of entry restores a mark from one to zero.
This proves injectivity on the full partial domain, not just on the
initialized tape.

\begin{theorem}[A delayed head move on an invariant plane]
\label{thm:sa-marked-head-plane}
The marked-head compiler gives an effective smooth mean-zero general
force on the unit torus, with all mixed derivatives bounded and period
one for $t\ge1$.  Its unique classical periodic zero-data solution is
globally smooth with bounded energy.  The particle from
$(1/4,1/4,1/2)$ enters $(5/8,7/8)^2\times\mathbb T$ exactly when the
original machine halts.  Only its loader depends on the input.
\end{theorem}
\begin{proof}
Put $N=(1/8,3/8)^2$, $J=(5/8,7/8)^2$.  Place disjoint rational state
squares in $N$ for nonhalt controls and in $J$ for halt controls.
For $k$ controls assigned to $(a,a+g)^2$, centers
$(a+ig/(k+1),a+g/2)$ and half-side $g/[4(k+1)]$ suffice.
Let $P_s$ map $[0,1]^2$ to the state square.  With base
$B=2|\Sigma|+1$, odd digits and $h_\sigma(x)=(d(\sigma)+x)/B$, an
instruction $(s,\sigma)\mapsto(s',\tau,d)$ has normalized maps
\[
\begin{array}{ll}
 d=1:&(x,h_\sigma(y))\mapsto(h_\tau(x),y),\\
 d=0:&(x,h_\sigma(y))\mapsto(x,h_\tau(y)),\\
 d=-1:&(h_\lambda(x),h_\sigma(y))\mapsto
          (x,h_\lambda(h_\tau(y))),\quad\lambda\in\Sigma.
\end{array}
\]
All free coordinates range over $[0,1]$.  Determinism and read symbols
separate sources.  At a common target control the incoming move is
fixed; distinct written symbols separate targets in the left first
digit for $d=1$, the right first digit for $d=0$, and the right pair
$(\lambda,\tau)$ for $d=-1$.  Hence the full target rectangles are
separated as well.  Composing with $P_s,P_{s'}$ preserves positivity
of the diagonal maps even when state squares have different sizes.

Apply Lemma~\ref{lem:sa-planar-routing} in $(0,1)^2$ with safe region
$N$.  To retain the linear-width alternative explicitly, shrink to a
common small rational half-side $r$, use sequential parking routes,
and linearly interpolate centers and widths with flat switches.  If
$0<\beta<1$ makes the rectangles enlarged by $1+\beta$ pairwise
disjoint within each endpoint family, the cutoff
\[
 \prod_{a=1}^2\left[1-\sigma\!\left(
 \frac{((x_a-c_a)/v_a)^2-1}{(1+\beta)^2-1}\right)\right]
\]
multiplied by $\dot c+\diag(\dot v_a/v_a)(x-c)$ gives the exact closed
rectangle path $c(t)+\diag(v_a(t)/v_a(0))(x_0-c(0))$.
Taking $r$ below one sixteenth of every relevant rational squared
clearance and endpoint half-side makes the small moving supports
miss the stationary ones.  Nonhalt paths remain in $N$ throughout.
This formula proves exactness on closed rectangles; the padded version
of Lemma~\ref{lem:sa-planar-routing} separately gives exterior neighborhoods.

Load the rational initial code from $(1/4,1/4)$ by the same planar
routing of a small square, confined to $N$ when the initial state is
nonhalting.  Follow by the unit step pulse.  Lift with
$h(z)=(z-1/2)\chi(z)$, where $\chi=1$ for $|z-1/2|\le1/8$ and is
supported in $|z-1/2|<1/4$.  Lemmas~\ref{lem:sa-invariant-lift} and
\ref{lem:sa-realization} give the mean-zero velocity and force with
pressure zero.  At unit samples after loading the particle is the
successive code.  Halt codes lie in $J$; a nonhalting run remains in
$N$ throughout loading and all later motions.  This proves the event.
The log's leftmost nonblank cell remains absolute cell one, so the
encoded tape also determines absolute head location.
\end{proof}

\subsection{A pointer that edits the neighboring cell}
\label{subsec:sa-neighbor-recorder}
The previous compiler advances the frontier through a separate control.
The next one advances it by editing the next cell directly.  This saves
a state, but its returning scan needs a different guard.  We state that
guard on the entire partial domain before considering initialized tapes.

Let $A$ be the work alphabet and $D=\{-1,0,1\}$.  Put
\[
 \mathcal R=(Q\setminus H)\times D\times A,\quad
 K=\{\bot,\mathsf{top}\}\sqcup\{[r]:r\in\mathcal R\},\quad
 \Gamma=A\times K\times\{0,1\}.
\]
Let $\mathcal S$ consist of the controls $C_{q,e},I_{q,e},O_r$,
with halting controls
$\mathcal H=\{C_{q,e}:q\in H\}$.  Index a configuration's tape relative
to its head by $w_i=(a_i,k_i,m_i)$.  A move $\ell$ after edits means
$w'_i=\widetilde w_{i+\ell}$.  Unmentioned components remain unchanged.
There are exactly five kinds of rules:
\begin{enumerate}
\item At $C_{q,e}$, $q\notin H$, $m_0=0$, let $r=(q,e,a_0)$ and
$\delta(q,a_0)=(q',b,d)$.  Write $a_0=b,m_0=1$, enter $O_r$, move $1$.
\item At $O_r$, $m_0=0$, $k_0\ne\mathsf{top}$, make no edits and move
$1$, retaining the control.
\item At $O_r$, $m_0=0$, $k_0=\mathsf{top}$, $k_1=\bot$, write
$k_0=[r],k_1=\mathsf{top}$, enter $I_{q',d}$, and move $-1$.
Here $q',d$ are determined by $r$.
\item At $I_{q',d}$, $m_0=0$, $k_1\ne\mathsf{top}$, make no edits
and move $-1$, retaining the control.
\item At $I_{q',d}$, $m_0=1$, clear that mark, enter $C_{q',d}$,
and move $d$.
\end{enumerate}

\begin{lemma}[Neighboring-cell recorder]\label{lem:sa-neighbor-compiler}
This partial map $G$ is injective.  Initialize it at $C_{q_0,0}$ with
head zero, finitely specified work tape, marks zero, and blank history
except for its pointer at absolute cell $a\in\{1,2\}$.
At its $n$th compute-state visit the work configuration is the machine's
$n$th configuration, marks are zero, records occupy $a,\ldots,a+n-1$,
and the pointer is at $T_n=a+n$.  If $h_n$ is the head location, the
next machine step takes exactly $2(T_n-h_n)+1$ applications of $G$.
In particular, its initialized run reaches $\mathcal H$ exactly when
the machine halts.
\end{lemma}
\begin{proof}
An output in $O_r$ came by a right move.  Its mark at relative position
$-1$ distinguishes the compute entry (one) from a loop (zero).  In the
first case $r$ restores the old state, incoming tag, and work symbol;
in the second there were no edits.  An output in $I_{q',d}$ came by a
left move.  Its history symbol at relative position $2$ is a pointer
precisely for the entry from $O_r$; the return-loop guard excludes that
case.  For this entry the record at relative position $1$ identifies
$r$ and restores the old pair $(\mathsf{top},\bot)$.  An output in
$C_{q',d}$ has a unique predecessor type: undo move $d$ and restore
mark one.  This proves injectivity on every actual rule output; it does
not assert surjectivity or impose initialized-tape restrictions.

Since $|h_n|\le n$, $T_n>h_n$.  The compute rule marks $h_n$; the
outgoing scan reaches $T_n$, records the old data, advances the pointer,
and returns from $T_n-1$.  Every return-scan cell is unmarked until the
marked one, and its right neighbor is never the new pointer.  Thus each
guard holds.  Clearing the mark and moving by $d$ gives the next machine
configuration, with $h_n+d\le T_n<T_n+1$.  The count is
$1+(T_n-h_n-1)+1+(T_n-h_n-1)+1$.
It is at most $4n+5$ for $a=2$ and at most $4n+3$ for $a=1$.
This proves the induction and includes initial halting.
\end{proof}

The two initial pointer positions remain distinct: the $a=2$
initialization uses two more local steps per ordinary instruction than
$a=1$, by the exact count in Lemma~\ref{lem:sa-neighbor-compiler}.

To encode $G$, enumerate every applicable control and full triple
$(v_{-1},v_0,v_1)$.  These triples include every tested and edited cell
and partition the domain into finitely many cylinders.  The left and
right source prefixes have lengths $1$ and $2$.  If the move is $\ell$
and the edited triple is $\widetilde v$, the target prefixes are
\begin{equation}\label{eq:sa-neighbor-prefixes}
 L'=(\widetilde v_{\ell-j})_{1\le j\le\ell+1},\qquad
 R'=(\widetilde v_{\ell+j-1})_{1\le j\le2-\ell}.
\end{equation}
The old free tails at indices $-2,-3,\ldots$ and $2,3,\ldots$ remain
unchanged and ordered.  Thus each image is a full cylinder; distinct
images are disjoint by injectivity.  Lemma~\ref{lem:sa-radix} consequently
gives full separated source and target rectangles, with linear parts
$\diag(B^{-\ell},B^\ell)$ when state-square scales agree.

\begin{theorem}[Two invariant-plane realizations]
\label{thm:sa-neighbor-realizations}
The neighboring-cell recorder has the following effective zero-data
forced Navier--Stokes realizations at fixed positive computable viscosity.
Both have all mixed derivatives bounded, period one for $t\ge1$, and
uniformly bounded kinetic energy.
\begin{enumerate}
\item On the unit torus, any tape differing from blank on a specified
finite set of cells is allowed.  The force is mean zero and general,
and the solution is unique in the classical periodic class.  The particle
from $(1/16,1/16,0)$ enters $\{x:1/2<x_1<7/8\}$ exactly when the
machine halts.
\item On $\mathbb R^3$, for a finite word starting at cell zero, the
force is general with one compact spatial support uniform in time.
The particle from the origin has negative first coordinate at some time
exactly when the machine halts.  Uniqueness holds for smooth solutions
with bounded velocity on finite intervals and
$u\in C_tH^2\cap C_t^1L^2$, $p\in C_tH^1$; no bound on the competing
velocity gradient is required.
\end{enumerate}
The repeated field depends only on the table; only the loader uses the
input.  Neither force is required to decay or be solenoidal.  The same
formulas are effective relative to arbitrary fixed real positive viscosity.
\end{theorem}
\begin{proof}
For the torus use $a=2$, odd digits, base $B=2|\Gamma|+1$, and
$\lambda=1/(16|\mathcal S|)$.  Enumerating each type from zero, assign
corners
\[
 b_s=(1/8+2j\lambda,1/4)\quad(s\notin\mathcal H),\qquad
 b_s=(5/8+2j\lambda,1/4)\quad(s\in\mathcal H).
\]
Use the code $b_s+\lambda(E(L),E(R))$.
The nonhalt squares lie in $0<x<1/4$ and the halt squares in
$1/2<x<7/8$.  For Euclidean space use $a=1$, base $B=2|\Gamma|$ and
digits $0,2,\ldots,B-2$.  Give halt controls offsets $-2,-4,\ldots$
and all others offsets $2,4,\ldots$, writing $H(s)$ for the offset.
Use $(H(s)+E(L),E(R))$.  Zero digits can put codes on rectangle boundaries;
our maps are exact on the full closed rectangles.  Nonhalt rectangles
lie in $x\ge2$, and halt codes have $x<0$.
In each construction the prefixes in \eqref{eq:sa-neighbor-prefixes}
give the separated rectangles required for planar routing.

Here are explicit choices retaining the two all-time avoidance bounds.
Let $N$ be the number of rectangle pairs.  In the torus chart park at
\[
 p_i=(x_*,i/(N+1)),\qquad
 x_*={\tfrac12}\min(\{1/4\}\cup\{c_{i1},d_{i1}\}),
\]
where $c_i,d_i$ are endpoint centers.  A bend with $0<x<1/4$ can be
chosen off the forbidden lines as follows: choose its abscissa below
all vertical forbidden abscissae and $1/4$, then its positive ordinate
below $1$ and below all positive values of the other lines there.
All segments stay in the chart; when both endpoints are nonhalt their
centers stay below $1/4$.  Let $\rho$ be the minimum of $1$, endpoint
half-sides, boundary clearances, and squared distances from each segment
to stationary centers.  Half-side $\epsilon=\rho/10$ gives disjoint
translation supports because $3\sqrt2\epsilon<\rho\le\sqrt\rho$.
In Euclidean space park on the horizontal axis to the right of every
endpoint and of $2$.  For a transfer use a bend $(m,m^2)$, where the
integer $m\ge2$ exceeds the Cauchy root bounds of all polynomials obtained
by restricting forbidden lines to that parabola.  Each polynomial is
nonzero of degree at most two, so this terminates.  A rational
$\epsilon<1$, below every endpoint half-side and one tenth of the
minimum squared segment clearance, gives disjoint supports.  A branch
joining nonhalt rectangles has centers at least $2$ throughout.

In either routing, shrink linearly to the small squares, translate them
in sequence, and expand linearly in the target rectangles, using the
closed-rectangle cutoff implementation in
Lemma~\ref{lem:sa-planar-routing}.  Put all slots inside $[1/4,3/4]$.
The normalized affine endpoint maps are exact, and the above clearances
show that every nonhalt torus path has $x<1/2$, while every nonhalt
Euclidean path has $x>0$.  The empty list uses zero.

Let $E_0,y_*$ be the rational initial planar codes.  On the torus load
from $a=(1/16,1/16)$ by the spatially constant field
$\dot\eta(t)(E_0-a)$, with $\eta$ flat from zero to one in
$[1/4,3/4]$.  Follow by the repeated rule pulse and lift with
$h(z)=\sin(2\pi z)/(2\pi)$.  In Euclidean space load the origin along
$\theta(t)y_*$ with a compact moving cutoff equal to one on a square
of half-side one and supported in its enlargement by one.  Follow by
the repeated planar pulse and lift with $h(z)=z\chi(z)$, where
$\chi=1$ on $[-1,1]$ and is supported in $(-2,2)$.
Lemma~\ref{lem:sa-invariant-lift} gives the exact planar motions and
Lemma~\ref{lem:sa-realization} gives the forces.  Its bounded-velocity
comparison refinement is exactly the second row's uniqueness class.

At time one both particles are their initial codes.  Subsequent rule
samples encode the run.  In the torus case the $n$th compute visit is
at $1+\sum_{j<n}(2(T_j-h_j)+1)$ with $T_j=j+2$; the Euclidean case
has $T_j=j+1$.  A halt code satisfies the observer, including an initial
halt after loading.  In a nonhalting run the torus loader stays below
$1/4$, and the Euclidean loader has nonnegative first coordinate; all
later paths obey the strict bounds proved above.  The history block's
known absolute starting cell $a$ recovers absolute head location as
well as the head-relative work tape.  All support and derivative bounds
come from finitely many prescribed pulses.
\end{proof}

\section{Marked tape words and a separate history stack}
\label{sec:sa-marked-five-stage}
The tape-track recorders retain the indexed work tape while a head
travels to the history frontier.  Here the entire finite work interval is
stored in the first stack and history in the second.  To reach the marked
head, the processor transfers the preceding work prefix onto the history
stack, performs the local update, and restores that prefix.  Its geometric
realization separates horizontal scaling into two successive operations,
making the intermediate observer bound explicit.
Normalize to one halt state $h$ and denote the work alphabet by $\Gamma$.

For each nonhalt instruction $(q,a)\mapsto(q',b,d)$ make one case $i$
if $d=0$, and, for each of $d=\pm1$, one case for each neighbor letter
$c\in\Gamma$ and one end case.  Write $q^+(i)=q'$ for the
destination state of case $i$.  Use alphabet
\[
 \mathcal A=\Gamma\sqcup\{\bar a:a\in\Gamma\}
       \sqcup\{E,\#\}\sqcup\{J_i\}_i\sqcup\{J_*\},
\]
where bars mark the head and $E$ pads the finite tape word.
The controls are $P_q,S_q$ (the latter only for nonhalt $q$), and $R_i$.

At a $P_q$ visit the first stack is a finite tape interval in its natural
left-to-right order, with exactly one marked head, followed by $E^\infty$.
The second is finite history followed by $E^\infty$.  Initially the
interval is cells $0,\ldots,\max(0,|w|-1)$, with its first symbol marked;
empty input gives the one-letter word $\bar\square$.  The initial
control is $P_{q_0}$ and the second stack is $E^\infty$.
The $S_q$ scan transfers the word left of the mark, the update writes and
moves the head mark (adjoining a blank precisely at an endpoint), and the
$R_i$ scan restores the left part.  Replacing $\#$ by $J_i$ finishes
the step.  If the left-of-head lengths before and after are $L_-,L_+$,
its count is $3+L_-+L_+$.  History determines the absolute left endpoint
as minus the number of left-end extensions.  At $P_h$ the halt loop
merely adds $J_*$ forever.

Entry and scan rules are
\[
 (P_q;\epsilon,\epsilon)\to(S_q;\epsilon,\#)\quad(q\ne h),\qquad
 (P_h;\epsilon,\epsilon)\to(P_h;\epsilon,J_*),
\]
\[
 (S_q;c,\epsilon)\to(S_q;\epsilon,c)\quad(c\in\Gamma).
\]
For the cases of $(q,a)\mapsto(q',b,d)$ the updates are
\[
\begin{array}{ll}
 d=0:&(S_q;\bar a,\epsilon)\to(R_i;\bar b,\epsilon),\\
 d=1:&(S_q;\bar a c,\epsilon)\to(R_i;\bar c,b),\quad c\in\Gamma,\\
     &(S_q;\bar a E,\epsilon)\to(R_i;\bar\square E,b),\\
 d=-1:&(S_q;\bar a,c)\to(R_i;\bar c b,\epsilon),\quad c\in\Gamma,\\
     &(S_q;\bar a,\#)\to(R_i;\bar\square b,\#).
\end{array}
\]
Complete them by
\[
 (R_i;\epsilon,c)\to(R_i;c,\epsilon)\quad(c\in\Gamma),\qquad
 (R_i;\epsilon,\#)\to(P_{q^+(i)};\epsilon,J_i).
\]
All tested and produced prefixes have length at most two.

For the full-domain check, ordinary versus marked first-stack tops
separate scans and updates at $S_q$; the displayed neighbor tests separate
the update cases.  Right tops separate the $R_i$ rules.  Into $S_q$,
entry writes $\#$ on the right while scans write ordinary letters.
Into $R_i$, its unique update has a marked first-stack top and return
loops have ordinary first tops.  Into $P_q$, distinct $J_i$ distinguish
case exits and $J_*$ distinguishes the halt loop.  Thus full source and
image cylinders are separately disjoint, and the initialized machine
simulation is faithful indefinitely.

\begin{theorem}[Marked words and a fixed middle arc]
\label{thm:sa-marked-five-stage}
This recorder gives an effective smooth mean-zero general force on the
unit torus, with all mixed derivatives bounded and period one for $t\ge1$.
Its unique classical periodic zero-data solution is global with bounded
energy.  The particle from $(1/4,1/4,0)$ enters
$(1/2,3/4)\times\mathbb T^2$ exactly when the machine halts.
The repeated protocol uses only the table and initialization uses the
input.  Separately, one may choose either a solenoidal force with the
same eventual periodicity, or a mean-zero force satisfying
$\|f(t)\|_{C^m}=O((1+t)^{-1})$ for every spatial order $m$, with all
mixed derivatives bounded and hence temporal $L^2$ in every $C^m$ norm.
These two alternatives are separate prescriptions.
\end{theorem}
\begin{proof}
Use odd digits and base $K=2|\mathcal A|+1$.  If there are $N$ controls,
let $\sigma_0=1/[100(N+1)]$, number them by $j(s)$, and put
\[
 o_s=(x_s,1/4+3\sigma_0j(s)),\qquad
 x_s=\begin{cases}5/8,&s=P_h,\\1/4,&s\ne P_h.\end{cases}
\]
Encode by $(o_s+\sigma_0(E(A),E(B)),0)$.
State squares are separated vertically by $2\sigma_0$.  Within a state,
Lemma~\ref{lem:sa-radix} gives a source or target coordinate gap at least
$\sigma_0K^{-2}$.  Write $D_r,D_r'$ for the matched rectangles,
$c_r,c_r'$ for their centers, and
\[
 \lambda_r=K^{|\alpha|-|\alpha'|},\quad
 \mu_r=K^{|\beta|-|\beta'|},\quad
 F_r(y)=c_r'+\diag(\lambda_r,\mu_r)(y-c_r).
\]
All side lengths are at most $\sigma_0$.

We describe five divergence-free fields whose successive unit-mass pulses
realize $F_r$.  Let $\rho_D$ equal one on the $\delta/2$ enlargement of
$D$ and be supported in its $\delta$ enlargement, where
$\delta=\sigma_0/(10K^2)$.  Put
$A_D=\partial_x((x-c_1)\rho_D)$; it is one on $D$ and has zero integral.
For $M$ rules, numbered $j_r=0,\ldots,M-1$, assign
\[
 z_r=1/4+\frac{j_r+1}{2(M+1)},\qquad m_z=\frac1{16(M+1)}.
\]
Let $\theta_r(z)$ be one on $[z_r-m_z,z_r+m_z]$, supported in its
$m_z$ enlargement, and set $B_r=\partial_z((z-z_r)\theta_r)$.
The layer supports are disjoint, $B_r=1$ near $z_r$, and $\int B_r=0$.
Take $\eta_r$ to be the product of $\theta_r$ and cutoffs equal to one
on $[c_{ra}-\sigma_0,c_{ra}+\sigma_0]$ with margin $\sigma_0$,
for $a=1,2$.  Define
\[
 H_r^x=(\log\lambda_r)(x-c_{r1})(z-z_r)\eta_r,
 \qquad H_r^y=(\log\mu_r)(y-c_{r2})(z-z_r)\eta_r,
\]
\[
\begin{aligned}
 W_1&=(0,0,\sum_r z_rA_{D_r}),\\
 W_2&=\sum_r(\partial_zH_r^x,0,-\partial_xH_r^x),\\
 W_3&=\sum_r(0,\partial_zH_r^y,-\partial_yH_r^y),\\
 W_4&=\sum_r B_r(z)(c_r'-c_r,0),\\
 W_5&=(0,0,-\sum_r z_rA_{D_r'}).
\end{aligned}
\]
All functions are smooth torus functions by zero extension of their
compact chart factors.  Every field is divergence free: mixed partials
cancel in $W_2,W_3$, and the others are constant in their flow directions.
Their means vanish by the derivative formulas, including $\int B_r=0$.
For $1\le j\le5$ let $\beta_j$ be the derivative of a flat step from
zero to one on $[(2j-1)/12,2j/12]$, and put
$W(\tau)=\sum_j\beta_j(\tau)W_j$, periodically extended.
These disjoint pulses are nonnegative and have integral one.

Starting from $(y,0)$, $y\in D_r$, the first pulse raises it to
$(y,z_r)$.  On that layer the second field is
\[
 ((\log\lambda_r)(x-c_{r1}),0,-(\log\lambda_r)(z-z_r)),
\]
so it scales the first offset by $\lambda_r$ and keeps $z=z_r$.
At intermediate pulse mass $b\in[0,1]$ its factor is $\lambda_r^b$;
the offset is at most $\sigma_0/2$ because both endpoint widths are at
most $\sigma_0$.  Thus the candidate path stays on the asserted plateau.
The third pulse analogously scales the second offset by $\mu_r$, still
on its plateau.  The fourth translates by $c_r'-c_r$, since
$B_r(z_r)=1$ and other layers vanish.  Target separation makes the fifth
pulse lower the point back to zero.  Smooth ODE uniqueness therefore gives
$(y,0)\mapsto(F_r(y),0)$ on every full source rectangle.
For a branch whose two controls are nonhalting, the first coordinate
throughout lies in
\[
 [1/4-\sigma_0/2,1/4+3\sigma_0/2],
\]
both while scaling and on the subsequent translation.  It avoids the
observed arc.  On the halt loop $\lambda_r=1$ and $c_{r1}'=c_{r1}$,
so the first coordinate stays unchanged throughout.

Let $y_{\rm in}$ be the rational initial code.  During $[0,1]$ use
$\beta_0(t)\cos(2\pi z)(y_{\rm in}-(1/4,1/4),0)$, where
$\beta_0$ is a unit-mass flat pulse on $[1/4,3/4]$.  It is divergence
free and mean zero, and the distinguished particle travels on $z=0$
to its initial code.  Follow with $W(t-1)$.  The resulting velocity $V$
is smooth at the joins, starts from zero, and has bounded derivatives.
Lemma~\ref{lem:sa-realization} gives its residual force and pressure zero.
The exact rectangle maps give the rule-code induction at times $1+k$.
A halt lies in $[5/8,5/8+\sigma_0]\subset(1/2,3/4)$; a nonhalting run
has a safe loader and stays in the displayed band at all subsequent
times.  This proves the first force contract.

For the decaying choice use the onto clock $s=\log(1+t)$ and
$\widehat V(t)=aV(s)$, $a=(1+t)^{-1}$.  Its force is
\[
 \widehat f=a^2(V_s-V+(V\cdot\nabla)V)(s)-\nu a\Delta V(s).
\]
Bounded derivatives of $V$ give the asserted $C^m$ estimate; every time
derivative is a finite sum of bounded derivatives of $V$ times powers
of $a$, so mixed derivatives are bounded.  Integrating $(1+t)^{-2}$
gives temporal $L^2$.  Since $s$ covers all nonnegative internal time,
the same exact material event is retained.  This choice replaces the
periodic forcing; it does not impose decay on that periodic field.
For the separate solenoidal periodic choice, let
$\Delta\phi=\operatorname{div}((V\cdot\nabla)V)$, $\int\phi=0$.
Proposition~\ref{prop:projection-l2} gives
$f_{\rm sol}=\partial_tV+(V\cdot\nabla)V-\nabla\phi-\nu\Delta V$
and $p=-\phi$, preserving all periodic bounds and the velocity.
All cutoffs, logarithms of positive rational scale factors, and finite
differentiations are effective; none depend on a future executed branch.
\end{proof}
